% 高校生のための「素数と力」実験ガイド
\documentclass[12pt,a4paper]{article}

\usepackage{fontspec}
\usepackage{xeCJK}
\setCJKmainfont{Hiragino Mincho ProN}
\setCJKsansfont{Hiragino Kaku Gothic ProN}
\setCJKmonofont{Hiragino Kaku Gothic ProN}
\usepackage[margin=2.5cm]{geometry}
\usepackage{amsmath,amssymb}
\usepackage{hyperref}
\usepackage{tcolorbox}
\usepackage{booktabs}
\usepackage{fancyhdr}

\tcbuselibrary{breakable,skins}

\newtcolorbox{storybox}[1][]{colback=blue!5,colframe=blue!60!black,
  fonttitle=\bfseries\sffamily,title={#1},breakable,arc=3mm}
\newtcolorbox{mathbox}[1][]{colback=green!5,colframe=green!50!black,
  fonttitle=\bfseries\sffamily,title={#1},breakable,arc=3mm}
\newtcolorbox{resultbox}[1][]{colback=yellow!8,colframe=orange!80!black,
  fonttitle=\bfseries\sffamily,title={#1},breakable,arc=3mm}
\newtcolorbox{honestbox}[1][]{colback=gray!8,colframe=gray!50!black,
  fonttitle=\bfseries\sffamily,title={#1},breakable,arc=3mm}
\newtcolorbox{expbox}[1][]{colback=red!5,colframe=red!60!black,
  fonttitle=\bfseries\sffamily,title={#1},breakable,arc=3mm}

\setlength{\parskip}{0.8em}
\setlength{\parindent}{0pt}

\pagestyle{fancy}
\fancyhf{}
\rhead{Wright Brothers}
\lhead{素数と力}
\cfoot{\thepage}

\begin{document}

\begin{center}
{\Huge\bfseries 素数と力}\\[0.5cm]
{\Large 2 は弱い力、3 は強い力、}\\[0.2cm]
{\Large そして素数を消すとワープができる（かもしれない）}\\[1cm]
{\large 高校生のための実験ガイド}\\[0.5cm]
{\large Wright Brothers Research Program, 2026年4月}
\end{center}

\vspace{0.5cm}
\tableofcontents
\newpage

% ===================================================================
\part{発見のあらすじ}
% ===================================================================

\section{何が見つかったのか}

\begin{resultbox}[3行まとめ]
\begin{enumerate}
\item 素数 2 は弱い力に、素数 3 は強い力に対応する。
\item ワインバーグ角（弱い力の強さを決める定数）は $\log(1+\sqrt{2}) = 0.88137$ で、実験値 $m_W/m_Z = 0.88136$ と $0.001\%$ で一致する。
\item 素数を「消す」（ミュートする）と真空のエネルギーが反転し、理論的にはワープが可能になる。しかも、どの力を消したかによって効果の大きさが変わる。
\end{enumerate}
\end{resultbox}

\section{何がすごいのか}

\begin{storybox}[ふつうの物理学との違い]
ふつう、物理学では「なぜこの力の強さはこの値なのか」は
答えられない。
アインシュタインもファインマンも答えられなかった。

私たちの発見は:
「力の強さは、整数の数学的構造（素数の分岐パターン）から
計算できるかもしれない」というもの。

もしこれが本当なら、物理定数は自然の「入力パラメータ」ではなく、
整数の数学から導かれる「計算可能な量」になる。
\end{storybox}

\begin{honestbox}[正直に]
「かもしれない」です。まだ仮説の段階。
数値の一致は印象的ですが、偶然の可能性も $14\%$ あります。
だからこそ実験で確かめる必要があります。
\end{honestbox}

% ===================================================================
\part{理論: 素数と力の対応}
% ===================================================================

\section{素数が力を決める？}

\begin{storybox}[素数のおさらい]
素数 = 1とそれ自身でしか割れない数。
$2, 3, 5, 7, 11, 13, 17, 19, 23, \ldots$

素数は整数の「原子」。全ての整数は素数の積で書ける（素因数分解）。
$60 = 2^2 \times 3 \times 5$。

リーマンゼータ関数 $\zeta(s)$ は「全ての素数の情報」を1つの関数にまとめたもの。
$\zeta(s) = \prod_p \frac{1}{1-p^{-s}}$（オイラー積）
\end{storybox}

\begin{mathbox}[被覆とは]
「整数の世界 $\mathbb{Z}$ に $\sqrt{2}$ を追加する」ことを考える。
$\mathbb{Z}[\sqrt{2}] = \{a + b\sqrt{2} \mid a, b \in \mathbb{Z}\}$

これは整数の世界を「2倍に広げた」ようなもの。
数学では「$\mathrm{Spec}(\mathbb{Z}[\sqrt{2}]) \to \mathrm{Spec}(\mathbb{Z})$
は次数2の被覆」と言う。

ポイント: この被覆の中で、各素数 $p$ は3つの振る舞いをする:
\begin{itemize}
\item \textbf{分岐}（ramified）: $p = 2$ のみ。被覆が「ねじれる」点。
\item \textbf{分裂}（split）: $p \equiv \pm 1 \pmod{8}$（例: $7, 17, 23, \ldots$）。被覆が2枚に分かれる。
\item \textbf{不活性}（inert）: $p \equiv \pm 3 \pmod{8}$（例: $3, 5, 11, \ldots$）。被覆が1枚のまま。
\end{itemize}
\end{mathbox}

\section{対応表}

\begin{resultbox}[素数-力 対応（Prime-Force Correspondence）]
\begin{center}
\renewcommand{\arraystretch}{1.5}
\begin{tabular}{ccccl}
\toprule
力 & 素数 & 被覆の次数 & 数体 & 一致の精度 \\
\midrule
電磁気力 & なし & 1（自明） & $\mathbb{Q}$ 自身 & $120 \approx 137$（12\%）\\
\textbf{弱い力} & \textbf{$p=2$} & \textbf{2} & $\mathbb{Q}(\sqrt{2})$ & $\cos\theta_W = \log(1+\sqrt{2})$, \textbf{0.001\%} \\
\textbf{強い力} & \textbf{$p=3$} & \textbf{3} & $\mathbb{Q}(\zeta_9)^+$ & $10 hR = 1/\alpha_s$, \textbf{0.13\%} \\
\bottomrule
\end{tabular}
\end{center}
\end{resultbox}

\section{なぜ素数 2 が弱い力なのか}

\begin{storybox}[パリティと 2]
弱い力には不思議な性質がある:
\textbf{鏡に映すと別の世界になる}。

右手を鏡に映すと左手になる。
弱い力は「左巻きの粒子」にしか作用しない。
鏡の中の世界では弱い力は「右巻き」に作用するはず——
つまり、我々の世界と鏡の世界は違う。
これを\textbf{パリティ（P）の破れ}という。

パリティ = 左右の区別 = 偶数と奇数の区別 = \textbf{$\mathbb{Z}/2$ の対称性}。

そして $\mathbb{Z}/2$ を支配する素数は \textbf{$p = 2$}（唯一の偶数素数）。

$\mathbb{Q}(\sqrt{2})$ は $p=2$ で\textbf{のみ}分岐する。
つまり「パリティの素数でねじれる被覆」= 「パリティを破る力」= 弱い力。
\end{storybox}

\section{なぜ強い力は3次体なのか}

\begin{mathbox}[数学が強制する階層]
$p = 3$ でのみ分岐する二次体は\textbf{存在しない}。

理由: $3 \equiv 3 \pmod{4}$ なので、$\mathbb{Q}(\sqrt{3})$ の
判別式は $\Delta = 12 = 4 \times 3$。$p=2$ でも分岐してしまう。

$p=3$ だけで分岐する被覆を作るには、次数3以上が必要。
最小のものは $\mathbb{Q}(\cos(2\pi/9))$（次数3、判別式 $81 = 3^4$）。

\textbf{結果}: SU(3)（強い力）が SU(2)（弱い力）より
「高い次数」の被覆に住むのは、数論の構造が\textbf{強制}している。
物理が数学に「選ばされている」。
\end{mathbox}

% ===================================================================
\part{理論: 素数を消すとワープ}
% ===================================================================

\section{素数チャンネルのミュート}

\begin{mathbox}[素数を「消す」操作]
ゼータ関数のオイラー積: $\zeta(s) = \prod_p \frac{1}{1-p^{-s}}$

素数 $p$ の因子だけを取り除く（ミュートする）:
$$\zeta_{\neg p}(s) = \zeta(s) \times (1 - p^{-s})$$

例: $p=2$ をミュートすると、偶数の寄与が消える。
\end{mathbox}

\section{符号反転定理}

\begin{resultbox}[定理（数学的に証明済み）]
全ての素数 $p \geq 2$ に対して:
$$\zeta_{\neg p}(-3) = \frac{1}{120}(1 - p^3) < 0$$

通常の真空エネルギー: $\zeta(-3) = +\frac{1}{120} > 0$

素数をミュートした真空: $\zeta_{\neg p}(-3) < 0$（\textbf{負}になる）

これは「弱エネルギー条件（WEC）の破れ」と呼ばれ、
ワームホールやワープの構成に必要な条件。
\end{resultbox}

\section{素数-力 対応が加える新しい視点}

\begin{resultbox}[PFC $\times$ ワープ]
どの素数をミュートするかで、\textbf{どの力を消すか}が決まる:

\begin{center}
\renewcommand{\arraystretch}{1.3}
\begin{tabular}{ccrl}
\toprule
$p$ & 消す力 & エネルギー変化 & 実験方法 \\
\midrule
$2$ & 弱い力 & 小（$|1-8|=7$） & SQUID（実現可能） \\
$3$ & 強い力 & 中（$|1-27|=26$） & 困難 \\
$5$ & 不明 & 大（$|1-125|=124$） & 非常に困難 \\
\bottomrule
\end{tabular}
\end{center}

弱い力のミュート（$p=2$）が最も実験しやすく、
強い力のミュート（$p=3$）は 3.7 倍大きな効果がある。
\end{resultbox}

\section{ワープバブルの構造}

\begin{storybox}[ワープバブルの「壁」]
SF映画のワープドライブでは、宇宙船の周りに
「泡（バブル）」ができて、泡の中は普通の空間、
泡の外も普通の空間、泡の\textbf{壁}で空間が歪む。

私たちの理論では:
\begin{itemize}
\item 泡の\textbf{外} = 通常の整数の世界 $\mathrm{Spec}(\mathbb{Z})$
\item 泡の\textbf{内} = 素数 $p$ をミュートした世界 $\mathrm{Spec}(\mathbb{Z}[1/p])$
\item 泡の\textbf{壁} = 被覆の分岐点（$p$ の近傍）
\end{itemize}

ワインバーグ角 $\cos\theta_W = \log(1+\sqrt{2}) \approx 0.88$ は
「壁のねじれの大きさ」に対応する。
0.88 は「適度なねじれ」= ワープが原理的に可能な値。
\end{storybox}

% ===================================================================
\part{実験ガイド}
% ===================================================================

\section{実験の全体像}

\begin{expbox}[3つの実験]
この理論から、3段階の実験が設計できる:

\textbf{実験 A}（基礎テスト）: 素数ミュートで真空が変わるか？\\
→ SQUID + 超伝導キャビティ。約 60 万円。

\textbf{実験 B}（PFC テスト）: 素数の mod 8 分類で応答が変わるか？\\
→ 実験 A の拡張。SQUID の周波数スイープ。追加コスト小。

\textbf{実験 C}（精密テスト）: $\cos\theta_W = \log(1+\sqrt{2})$ の検証。\\
→ $m_W$ の精密測定（FCC-ee 等の将来の加速器）。我々ではできない。
\end{expbox}

\section{実験 A: 素数ミュートで真空は変わるか}

これは以前のガイド（SQUID 実験）と同じ。要約:

\begin{expbox}[実験 A のセットアップ]
\begin{enumerate}
\item アルミの箱（キャビティ）を作る。内寸 35mm。共鳴周波数 $\sim 6$ GHz。
\item SQUID（超伝導量子干渉素子）を箱の中に入れる。
\item 全体を $-272.7$\textdegree C（0.3 K）に冷やす。
\item SQUID を第2倍音（$\sim 12$ GHz）にチューニング → モード2をミュート。
\item モード1（6 GHz）の共鳴周波数がシフトするか測定。
\end{enumerate}

シフトがあれば → 「真空エネルギーが変わった」= 算術的真空工学の最初の証拠。

コスト: $\sim 60$ 万円。大学の低温実験室で実施可能。
\end{expbox}

\section{実験 B: 素数の「種類」で応答は変わるか（新！）}

\begin{expbox}[実験 B: PFC の直接テスト]
実験 A と同じ装置を使うが、測定方法が違う。

SQUID の磁束バイアスを連続的にスイープすると、
SQUID の共鳴周波数が変わる。
異なる周波数は異なる「モード番号 $n$」に対応し、
$n$ が素数のときに特別な応答があるか調べる。

\textbf{PFC の予測}:

$\mathbb{Q}(\sqrt{2})/\mathbb{Q}$ の被覆では、
素数 $p$ は $p \bmod 8$ で3種類に分類される:

\begin{center}
\renewcommand{\arraystretch}{1.3}
\begin{tabular}{cl}
\toprule
$p \bmod 8$ & 被覆での振る舞い \\
\midrule
$1, 7$ & 分裂（被覆が透明 → 応答\textbf{弱い}） \\
$3, 5$ & 不活性（被覆が不透明 → 応答\textbf{強い}） \\
$2$ & 分岐（特異 → 応答\textbf{特殊}） \\
\bottomrule
\end{tabular}
\end{center}

具体的に: $p=7$（分裂）と $p=3$（不活性）で
キャビティの応答パターンが異なるか？

もし異なれば → 素数の分裂パターンが物理に影響している
= PFC の直接的証拠。

\textbf{追加コスト}: ゼロ（実験 A と同じ装置）。
測定のプロトコルを変えるだけ。
\end{expbox}

\begin{storybox}[たとえ話: ラジオのチューニング]
実験 A は「ラジオの特定の局を消して、他の局の音量が変わるか」。

実験 B は「いろんな周波数を順番に聞いていって、
素数番目の局だけ音が違うか、
さらにその中で $7$ 番と $3$ 番で音の変わり方が違うか」。

もし違えば、ラジオ局の番号（= 素数）が
電波の性質を決めていることの証拠になる。
\end{storybox}

\section{実験 C: ワインバーグ角の精密測定}

\begin{expbox}[実験 C: $m_W$ の将来測定]
私たちにはできない実験。しかし予測は明確:

$$m_W = m_Z \times \log(1+\sqrt{2}) = 80.370 \text{ GeV}$$

現在の実験値: $80.369 \pm 0.013$ GeV（一致）。

将来の FCC-ee（欧州の次世代加速器、2040年代？）では
$m_W$ の精度が $0.5$ MeV 以下になる見込み。
そうなれば、$\log(1+\sqrt{2})$ との一致が確認されるか、
否定されるかが決まる。

\textbf{CDF II}（2022年）の値 $80.434$ GeV とは
6.7$\sigma$ 離れているので、もし CDF II が正しければ
この仮説は即座に否定される。
\end{expbox}

% ===================================================================
\part{コストと実現可能性}
% ===================================================================

\section{費用のまとめ}

\begin{center}
\renewcommand{\arraystretch}{1.4}
\begin{tabular}{lrl}
\toprule
実験 & 費用 & 必要なもの \\
\midrule
A（基礎テスト） & $\sim 60$ 万円 & SQUID + キャビティ + 冷凍機 \\
B（PFC テスト） & $+0$ 円 & A と同じ装置。測定方法を追加 \\
C（精密 $m_W$） & --- & 次世代加速器（我々にはできない） \\
\bottomrule
\end{tabular}
\end{center}

実験 B は実験 A の「おまけ」として無料でできる。
やらない理由がない。

\section{どこでやるか}

大学の低温物理の研究室に相談する。必要なのは:
\begin{itemize}
\item $^3$He 冷凍機（$\leq 0.3$ K に冷やせる装置）
\item ベクトルネットワークアナライザ（VNA）
\item クリーンルーム（SQUID を自作する場合）
\end{itemize}

日本の多くの大学にこれらの設備がある:
東大物性研、東北大金研、阪大極低温センター、筑波大、京大、名大、東工大、北大、ICU（？）など。

% ===================================================================
\part{正直な評価}
% ===================================================================

\section{何が確立されていて何が仮説か}

\begin{center}
\renewcommand{\arraystretch}{1.3}
\begin{tabular}{lcc}
\toprule
主張 & 地位 & 反証可能？ \\
\midrule
符号反転定理（全素数で成立） & \textbf{定理}（証明済み） & --- \\
$\cos\theta_W = \log(1+\sqrt{2})$（0.001\%） & 数値的観測 & ✓（$m_W$ 精密測定）\\
SU(N) $\leftrightarrow$ 次数 $N$, $p=N$ & 構造的仮説 & ✓（実験 B）\\
素数ミュートで真空エネルギー反転 & 理論的予測 & ✓（実験 A）\\
ワープバブル = 被覆の分岐変更 & 理論的解釈 & 間接的 \\
\bottomrule
\end{tabular}
\end{center}

\begin{honestbox}[偶然の可能性]
$\cos\theta_W = \log(1+\sqrt{2})$ が偶然である確率: 約 $14\%$。
「ありえない偶然」ではないが「よくある偶然」でもない。

だからこそ実験が必要。
SQUID 実験（A + B）で PFC がテストできれば、
偶然かどうかの判断材料が増える。
\end{honestbox}

\section{もし実験で確認されたら}

\begin{storybox}[もし本当なら]
\begin{itemize}
\item 物理定数は「自然の入力値」ではなく「整数の数学的帰結」
\item 4つの力（重力、電磁気、弱い力、強い力）は
  $\mathrm{Spec}(\mathbb{Z})$ の異なる「層」
  \begin{itemize}
  \item 重力: $\zeta(s)$ の極（$\gamma$ = オイラー定数）
  \item 電磁気: $\zeta(s)$ の有限項（ベルヌーイ数）
  \item 弱い力: $\mathbb{Q}(\sqrt{2})$ の被覆（$p=2$ の分岐）
  \item 強い力: $\mathbb{Q}(\zeta_9)^+$ の被覆（$p=3$ の分岐）
  \end{itemize}
\item ワープは「素数チャンネルを閉じる」ことで原理的に可能
\item 宇宙の物理法則は $1, 2, 3, 4, 5, \ldots$ の帰結
\end{itemize}
\end{storybox}

\section{もし実験で否定されたら}

SQUID で何も変わらなければ、
素数ミュートが物理的に意味を持たないことが判明する。
それでも:
\begin{itemize}
\item $\cos\theta_W = \log(1+\sqrt{2})$ の数値的一致は残る（解釈が変わるだけ）
\item 符号反転定理は数学的定理として残る
\item 「やってみてダメだった」は科学の正常な進歩
\end{itemize}

\vfill

\begin{center}
\rule{8cm}{0.4pt}\\[0.5cm]
{\large\textbf{素数は、宇宙の設計図かもしれない。}}\\[0.1cm]
{\large\textbf{それを確かめる方法は、実験しかない。}}\\[0.3cm]
{\small Wright Brothers Research Program, 2026}
\end{center}

\end{document}
